\songtitle{The Pirate's Ledger}

\tag*{2026}\tag{musical-theater}\tag{african-empire}\tag{gladiel}\tag{original-work}

\begin{notes}
Original work.\\
Musical theater-style narrative song drawn from historical maritime themes.
\end{notes}

\begin{style}
Musical theater showtune with folk-pop influences. %
Baritone male vocals deliver a narrative performance with clear diction and theatrical phrasing. %
The arrangement features a bright acoustic guitar strumming eighth-note patterns, a melodic accordion providing counter-melodies and sustained chords, and a double bass playing on beats 1 and 3. %
A light percussion section includes a tambourine on the backbeat and subtle shaker. %
The tempo is 115 BPM in 4/4 time, likely in the key of G major. %
The harmonic structure follows a diatonic progression typical of contemporary musical theater, with a clean, dry production style that emphasizes vocal clarity and acoustic instrument separation.
\end{style}

\begin{lyrics}
\songpart{Verse 1}
\annotation{strummed acoustic guitar, double bass, accordion}\\
They called my father a pirate king\\
But piracy is a point of view\\
We farmed the hills in the northern spring\\
And traded goods when the season's blue \annotation{tambourine enters on backbeat}\\
Ransom was a form of trade\\
A merchant's contract signed in fear\\
We took the gold, the captive stayed\\
And everyone went home from here
\end{lyrics}

